\documentclass[12pt,a4paper,oneside]{article} %стиль оформления - статья с параметрами: шрифт 12, формат бумаги А4, печать с одинаковыпи полями на всех страницах - oneside
\usepackage[T2A]{fontenc} %поддержка кириллицы
\usepackage[utf8]{inputenc} %кодировка символов
\usepackage[english,russian]{babel} %пакет поддержки русского языка с элементами английского
\frenchspacing
\usepackage[dvips]{graphicx} %разрешить включение PostScript-графики
\graphicspath{{images/}}
\usepackage{setspace} %"умное" расстояние между строк
\usepackage[usenames,dvipsnames]{color}
\usepackage{hhline} %работа с таблицами
\usepackage[12pt]{extsizes} %размер шрифта
\usepackage{fancyhdr} %установка колонтитулов

\usepackage{geometry} %ручная установка полей
\geometry{top=1cm} %поле сверху
\geometry{bottom=1cm} %поле снизу
\geometry{left=2cm} %поле слева
\geometry{right=2cm} %поле справа

\begin{document}
\singlespacing % межстрочный одиночный интервал для всего текста
\pagestyle{empty} % нумерация выкл.

%Н:задание на уир и кп
\newpage
\begin{center}
МИНИСТЕРСТВО ОБРАЗОВАНИЯ И НАУКИ РОССИЙСКОЙ ФЕДЕРАЦИИ

Федеральное государственное бюджетное образовательное учреждение \linebreak
высшего профессионального образования

\textbf{"Национальный исследовательский ядерный университет "МИФИ"}
\\[10pt]

\begin{figure}[h]
\includegraphics[trim=-10 40 30 0,width=0.25\linewidth]{mifi}
\end{figure}

ФАКУЛЬТЕТ КИБЕРНЕТИКИ
\\[15pt]

КАФЕДРА ИНФОРМАЦИОННЫЕ ТЕХНОЛОГИИ
\\[50pt]

\huge{\textbf{Задание на УИР и КП}}
\\[20pt]

\begin{flushleft}
\large{Студентке гр.\underline{К7-361 Лаврентьевой Марине Олеговне\hspace{5,2cm}}}
\\[40pt]
\end{flushleft}

\Large{\textbf{ТЕМА УИР и КП}}

\begin{flushleft}
\large{\underline{Разработка модуля преобразования геоданных из формата\hspace{4cm}} \linebreak
\underline{OpenStreetMap в формат OpenGIS\hspace{9,4cm}}}
\\[60pt]
\end{flushleft}

\large{\textbf{ЗАДАНИЕ}}
\end{center}

\large{\begin{enumerate}
\item Изучение структуры xml-документа с геоданными в формате \linebreak
OpenStreetMap.
\item Построение XPath-запросов для извлечения геоданных об объектах \linebreak
заданного типа.
\item Изучение стандарта OpenGIS и разработка структуры базы \linebreak
данных для хранения геоданных в формате OpenGIS.
\item Построение SQL-запросов для сохранения и извлечения геоданных \linebreak
в формате OpenGIS.
\item Реализация модуля преобразования геоданных из формата \linebreak
OpenStreetMap в формат OpenGIS.
\end{enumerate}}

\newpage
\normalsize{\begin{flushleft}
Место выполнения УИР и КП \underline{\hspace{11,15cm}}

Руководитель \underline{\hspace{10,95cm} (Муравьёв С.К.)}
\\[8pt]
Дата выдачи задания \underline{\hspace{12,8cm}}
\end{flushleft}

\newpage
\subsubsection*{ОТЗЫВ О РАБОТЕ СТУДЕНТА}

\vfill
Руководитель \underline{\hspace{8cm}\glqq{\hspace{1cm}}\grqq\hspace{4cm}2011г.}}
%К: задание на уир и кп
\end{document}
